\section{The limiting configuration theorem}
\label{sec:config}

\subsection{The Gram matrix}

We consider the rescaled observations
\begin{equation}
u_{ij}=\frac{x_{ij}-(\mu_1+\mu_2)/2}{\sqrt{\theta}}
=\eps_i\frac{\mu}{2\sqrt{\theta}}+\frac{v_{ij}+\xi_{ij}}{\sqrt{\theta}}.
\label{eq:4.1}
\end{equation}
Note that the translation does not change the weight vector of the SVM and the
scaling changes $\|w\|$ and $b$ simultaneously, so that the classification result
is not affected. We define the Gram matrix by
\begin{equation}
G=\bigl(u_{ij}^T u_{kl}\bigr)_{(i,j),(k,l)\in\cI}\in\R^{N\times N},
\label{eq:4.2}
\end{equation}
where $(i,j)$ gives the row and $(k,l)$ gives the column according to the order
\eqref{eq:2.2}. We also define
\[
D=\diag(\underbrace{c_{10},\dots,c_{10}}_{n_1},
\underbrace{c_{20},\dots,c_{20}}_{n_2}),
\]
whose components are arranged in the same order.

\subsection{The main theorem}

\begin{theorem}
\label{thm:1}
Assume \textup{(A1)} to \textup{(A4)} and \textup{(S)}. Then, as $d\to\infty$,
it holds that
\[
G\wto G^*=G_0^*+D,
\]
where the $((i,j),(k,l))$ component of $G_0^*$ is given by
\begin{equation}
\begin{aligned}
g^*_{(i,j),(k,l)}
&=\frac{\eps_i\eps_k\delta}{4}
+\frac{\eps_i\sqrt{\delta}}{2}\sum_{t\le m_k}\sqrt{c_{kt}}\,\beta_{kt}z_{klt}
+\frac{\eps_k\sqrt{\delta}}{2}\sum_{s\le m_i}\sqrt{c_{is}}\,\beta_{is}z_{ijs}\\
&\qquad+\sum_{s\le m_i}\sum_{t\le m_k}\sqrt{c_{is}c_{kt}}\,r_{st}^{(ik)}z_{ijs}z_{klt},
\end{aligned}
\label{eq:4.3}
\end{equation}
with $r_{st}^{(11)}=r_{st}^{(22)}=\delta_{st}$, $r_{st}^{(12)}=r_{st}$ and
$r_{st}^{(21)}=r_{ts}$. Moreover, when $c_{10}>0$ and $c_{20}>0$, it holds that
$G^*\succ O$ for every realization of $Z$.
\end{theorem}

\begin{proof}
(i) Convergence of the components. From~\eqref{eq:4.1} we have
\[
u_{ij}^Tu_{kl}
=\underbrace{\frac{\eps_i\eps_k\Delta}{4\theta}}_{\mathrm{(I)}}
+\underbrace{\frac{\eps_i}{2}\frac{\mu^T(x_{kl}-\mu_k)}{\theta}}_{\mathrm{(II)}}
+\underbrace{\frac{\eps_k}{2}\frac{\mu^T(x_{ij}-\mu_i)}{\theta}}_{\mathrm{(III)}}
+\underbrace{\frac{(x_{ij}-\mu_i)^T(x_{kl}-\mu_k)}{\theta}}_{\mathrm{(IV)}}.
\]
The term~(I) converges to $\eps_i\eps_k\delta/4$ from~(A3). The terms~(II)
and~(III) converge to the second and the third terms of~\eqref{eq:4.3} from
Lemma~\ref{lem:3}. As for~(IV), when $(i,j)\neq(k,l)$, it converges to the fourth
term of~\eqref{eq:4.3} from~\eqref{eq:3.1} for $i=k$ and from Lemma~\ref{lem:2} for
$i\neq k$. When $(i,j)=(k,l)$, it converges to the fourth term plus $c_{i0}$ from
\eqref{eq:3.2}. This gives the diagonal matrix $D$.

Since the number of the components is $N^2$, which does not depend on $d$, and all
of them are written as a common continuous function of $Z^{(d)}$ plus $o_P(1)$, the
convergence of the matrix follows from~(A4) together with the continuous mapping
theorem and Slutsky's theorem.

(ii) Positive definiteness. From Lemma~\ref{lem:6} in Section~\ref{sec:config-exist},
$G_0^*$ is the Gram matrix of a family of vectors $\tilde u_{ij}$ in $\R^K$, where
$K=1+m_1+m_2$. Hence, for any $a=(a_{ij})\in\R^N$, it holds that
\[
a^TG_0^*a=\Bigl\|\sum_{(i,j)\in\cI}a_{ij}\tilde u_{ij}\Bigr\|^2\ge 0.
\]
Note that this holds for every realization of $Z$, because
$(\tilde u_{ij})$ is a fixed family of vectors in $\R^K$ once the realization is
given. On the other hand, we have
\[
a^TDa\ge\min(c_{10},c_{20})\|a\|^2.
\]
Hence, when $c_{10}>0$ and $c_{20}>0$, for any $a\neq 0$ it holds that
$a^TG^*a>0$.
\end{proof}

\subsection{Existence of the limiting configuration}
\label{sec:config-exist}

In the proof of Theorem~\ref{thm:1}, we used the fact that $G_0^*$ is the Gram
matrix of a family of vectors. This fact is not trivial. If one specifies the
values of the inner products of $N$ vectors arbitrarily, there does not necessarily
exist a family of vectors which realizes them. For example, one cannot find three
unit vectors whose mutual inner products are all equal to $-1$. Hence we need to
prove the existence.

\begin{lemma}
\label{lem:6}
Assume \textup{(A3)} with $\Delta>0$ and let $K=1+m_1+m_2$. Then there exists a
family of non-random vectors $e_\mu$ and $\tilde h_{is}$ ($i=1,2$, $s\le m_i$) in
$\R^K$ such that
\begin{equation}
\|e_\mu\|^2=1,\quad
e_\mu^T\tilde h_{is}=\beta_{is},\quad
\tilde h_{is}^T\tilde h_{it}=\delta_{st},\quad
\tilde h_{1s}^T\tilde h_{2t}=r_{st}.
\label{eq:4.4}
\end{equation}
Moreover, putting
\begin{equation}
\tilde u_{ij}
=\eps_i\frac{\sqrt{\delta}}{2}\,e_\mu
+\sum_{s\le m_i}\sqrt{c_{is}}\,z_{ijs}\,\tilde h_{is},
\quad(i,j)\in\cI,
\label{eq:4.5}
\end{equation}
it holds that $\tilde u_{ij}^T\tilde u_{kl}=g^*_{(ij),(kl)}$.
\end{lemma}

\begin{proof}
We proceed in four steps.

\textit{Step~1.} We write the target inner products as a matrix. Let
\[
\cJ=\{\mu\}\cup\{(i,s):i=1,2,\,s\le m_i\},
\quad|\cJ|=K,
\]
and define a $K\times K$ symmetric matrix $\Gamma$ whose components are the
right-hand sides of~\eqref{eq:4.4}, that is,
$\Gamma_{\mu\mu}=1$, $\Gamma_{\mu,(i,s)}=\beta_{is}$,
$\Gamma_{(i,s),(i,t)}=\delta_{st}$, $\Gamma_{(1,s),(2,t)}=r_{st}$.
What we have to show is that there exists a family of vectors in $\R^K$ whose Gram
matrix is $\Gamma$.

\textit{Step~2.} We show that $\Gamma\succeq O$. We do not evaluate any inequality.
Instead, we carry the positive semi-definiteness which holds for finite $d$ to the
limit. For each $d$, we take the vectors
\[
e_\mu^{(d)}=\mu/\sqrt{\Delta},\qquad h_{is}\quad(i=1,2,\,s\le m_i),
\]
which actually exist in $\R^d$. Let $\Gamma^{(d)}$ be the Gram matrix of these $K$
vectors. Then $\Gamma^{(d)}\succeq O$ by definition, because for any
$a\in\R^K$ it holds that
\begin{equation}
a^T\Gamma^{(d)}a
=\Bigl\|a_\mu e_\mu^{(d)}+\sum_{i,s}a_{(i,s)}h_{is}\Bigr\|^2\ge 0.
\label{eq:4.6}
\end{equation}
The components of $\Gamma^{(d)}$ are given by
\[
(e_\mu^{(d)})^Te_\mu^{(d)}=1,\quad
(e_\mu^{(d)})^Th_{is}=b_{is},\quad
h_{is}^Th_{it}=\delta_{st},\quad
h_{1s}^Th_{2t}=\rho_{st}.
\]
Note that the first and the third equalities hold exactly for every $d$, because
$\|e_\mu^{(d)}\|=1$ and $H_i$ is an orthogonal matrix. As for the second and the
fourth, we have $b_{is}\to\beta_{is}$ and $\rho_{st}\to r_{st}$ from~(A3). Since
$K$ does not depend on $d$, the componentwise convergence gives
$\Gamma^{(d)}\to\Gamma$. Finally, fixing $a\in\R^K$ and letting $d\to\infty$ in
\eqref{eq:4.6}, we obtain $a^T\Gamma a\ge 0$ because the limit of a sequence of
non-negative numbers is non-negative. Since $a$ is arbitrary, we have
$\Gamma\succeq O$. This is nothing but a direct verification that the cone of
positive semi-definite matrices is closed.

\textit{Step~3.} We construct the vectors from $\Gamma\succeq O$. Let
$\Gamma=Q\Lambda Q^T$ be the spectral decomposition, where $Q$ is an orthogonal
matrix and $\Lambda=\diag(\ell_1,\dots,\ell_K)$ with $\ell_r\ge 0$. Let
$B=\Lambda^{1/2}Q^T$, which is well defined because $\ell_r\ge 0$, and put
$B^TB=\Gamma$. Then we define $e_\mu$ and $\tilde h_{is}$ as the columns of $B$
corresponding to the indices of $\cJ$. Since $B^TB$ is the inner product of the
columns of $B$, the equality $B^TB=\Gamma$ is exactly the four equalities in
\eqref{eq:4.4}. We note that $\Gamma$ is determined only by the non-random
quantities $\beta_{is}$ and $r_{st}$, so that $e_\mu$ and $\tilde h_{is}$ are
non-random.

\textit{Step~4.} We check the Gram matrix of~\eqref{eq:4.5}. From the bilinearity
we have
\begin{align*}
\tilde u_{ij}^T\tilde u_{kl}
&=\frac{\eps_i\eps_k\delta}{4}\|e_\mu\|^2
+\frac{\eps_i\sqrt{\delta}}{2}\sum_{t\le m_k}\sqrt{c_{kt}}\,z_{klt}\,e_\mu^T\tilde h_{kt}\\
&\quad+\frac{\eps_k\sqrt{\delta}}{2}\sum_{s\le m_i}\sqrt{c_{is}}\,z_{ijs}\,e_\mu^T\tilde h_{is}\\
&\quad+\sum_{s\le m_i}\sum_{t\le m_k}\sqrt{c_{is}c_{kt}}\,z_{ijs}z_{klt}\,
\tilde h_{is}^T\tilde h_{kt}.
\end{align*}
Substituting~\eqref{eq:4.4}, we obtain~\eqref{eq:4.3} term by term.
\end{proof}

\begin{remark}
\label{rem:2}
Lemma~\ref{lem:6} gives $G_0^*\succeq O$, but it does not give the positive
definiteness. In fact, in the setting of Section~\ref{sec:example} we have
$r_{11}=1$, that is, $\tilde h_{11}=\tilde h_{21}$, so that $\Gamma$ is singular.
In such a case, $G_0^*$ degenerates into a proper subspace of $\R^K$ and $G_0^*$ is
singular when $N>K$. Hence the positive definiteness of $G^*$ comes only from the
diagonal matrix $D$, that is, from the non-spiked noise. The condition $c_{i0}>0$
in~(S) is essential in this sense. When $c_{i0}=0$, the affine independence of the
data is lost and the arguments in Lemmas~\ref{lem:4} and~\ref{lem:5} break down.
\end{remark}

\begin{remark}
\label{rem:3}
When $\Delta=0$, the quantity $e_\mu$ is not defined. On the other hand, the
coefficient $\sqrt{\delta}/2$ of $e_\mu$ in~\eqref{eq:4.5} also vanishes because
$\delta=0$. Hence it suffices to remove the index $\mu$ from $\cJ$, put
$K=m_1+m_2$ and repeat the same argument. The first three terms of~\eqref{eq:4.3}
also vanish.
\end{remark}

\subsection{Discussion on the construction}

The construction in Lemma~\ref{lem:6} may look artificial. We give the idea behind
it, which is in fact simple.

First, we note that the SVM sees the data only through the inner products. As we
show in Section~\ref{sec:svm}, the dual problem involves only the Gram matrix $G$,
and the discriminant function involves only $u_p^Tu_0$. Hence the SVM is invariant
under congruent transformations, and we do not have to ask whether the points
themselves converge. We only have to ask whether the table of the inner products
converges. Lemmas~\ref{lem:1} to~\ref{lem:3} do exactly this, and the problem of
the diverging dimension is reduced to the convergence of an $N\times N$ matrix.

Second, however, we also want a geometric picture corresponding to the statement
that HDLSS data are located at the vertices of a regular simplex. This leads to
the inverse problem: does there exist a configuration of points which realizes the
limiting table of the inner products, and if so, in how many dimensions? This is
the problem of the classical multidimensional scaling.

Third, the tool for the inverse problem is the standard fact that, for a real
symmetric matrix $\Gamma$, the conditions $\Gamma\succeq O$, $\mathrm{rank}(\Gamma)=r$
for some $r$, and $\Gamma$ is a Gram matrix of some family of vectors, are
equivalent. Hence the proof of the positive semi-definiteness is itself the proof
of the existence of the configuration, and what we have to show is reduced to one
inequality.

Fourth, the positive semi-definiteness is obtained for free from the existence in
finite dimensions. For finite $d$, the vectors $e_\mu^{(d)}$ and $h_{is}$ actually
exist in $\R^d$, and the Gram matrix of vectors which actually exist is of course
positive semi-definite. We only have to carry this property to the limit, and the
condition~(A3) is assumed exactly for this purpose.

Finally, the dimension $K=1+m_1+m_2$ is the number of the directions which survive
in the limiting table of the inner products, namely, the mean difference direction
and the spiked directions. The non-spiked directions are killed by
Lemmas~\ref{lem:1} to~\ref{lem:3} except for the diagonal term $c_{i0}$.

We note that the choice of $B$ is not unique. One may use the Cholesky
decomposition or the reproducing kernel Hilbert space construction. We use the
spectral decomposition because it works without any modification even when $\Gamma$
is singular, as we remarked in Remark~\ref{rem:2}.

\begin{remark}[Collapse of the geometric representation]
\label{rem:4}
In the existing theory, which corresponds to $m_i=0$, the second to the fourth
terms of~\eqref{eq:4.3} vanish and $G^*$ is a deterministic matrix. This is the
geometric representation, and it is the reason why the SVM solution can be written
in a closed form. Under the spiked model, $G^*$ is a random matrix depending on
$Z$. That is, HDLSS data converge to a random configuration in a
$(1+m_1+m_2)$-dimensional space together with $N$ mutually orthogonal noise axes
whose lengths are $\sqrt{c_{i0}}$. All the arguments in the subsequent sections are
reduced to this finite-dimensional random limiting problem.
\end{remark}
